\section{Evaluate a causal coding workflow beyond its test-pass rate}
\label{exec19:evalcode-scope}

A causal-analysis assistant can generate syntactically correct code, execute a
query, and pass tests while answering the wrong banking question. A repayment
model might use postdecision information; an effect estimator might target the
wrong cohort; a reporting agent might convert an execution failure to a zero
effect. Evaluation must therefore identify the task, the observations on which
performance is measured, and the protocol that decides whether an output is
usable. Those choices connect the software interface to the scientific target
and to the economic case for a prototype.

Two selected historical surveys help organize this distinction. Chang and
coauthors' evaluation survey, inspected at its 6 July 2023 initial version,
separates what to evaluate, where to evaluate and how to evaluate. Its summary
explicitly limits its success and failure conclusions to particular datasets.
It discusses automated and human assessment, prompt robustness, complete
behavioral tests and evolving benchmarks. It also asks whether a dynamically
constructed test is truly out of distribution. Dynamic generation is thus a
design option, not a certificate of novelty or freedom from contamination.
These are historical organizing arguments, not a current ranking of models.
\citep{exec19_surveys_evaluation}

Jiang and coauthors' code-generation survey, inspected at its 1 June 2024
initial version, covers data curation, model adaptation, repository and agent
workflows, evaluation and applications. It distinguishes text/structural
similarity, execution metrics, human assessment and model judging. The inspected
performance tables compile cited model results; they are not a harmonized
experiment reproduced here. Its framework descriptions likewise need the
constituent technical originals before a package or claimed score can enter
our admitted design. The selected PDF's template footer says September 2024,
whereas the archive and first-page stamp identify the June initial version.
That footer does not establish journal publication or equivalence to the later
record that led to discovery.~\citep{exec19_surveys_code}

\subsection{A metric must define its sampling pool and selector}
\label{exec19:evalcode-selection}

The code survey's equation for pass-at-$k$ uses a finite-pool expression of the
form $1-\binom{n-c}{k}/\binom{n}{k}$. Its surrounding prose, however, describes
$c$ as the number of correct samples within the selected $k$, despite defining
$n$ as the larger generated pool. These definitions conflict. Before adopting
the expression, the implementation must fix which pool the count refers to.
The following derivation is our own finite counting argument, not an attribution
of an uninspected original theorem.~\citep{exec19_surveys_code}

Let a fixed pool contain $n=10$ candidates, of which exactly $c=3$ pass a
stipulated test suite. Choose $k=2$ candidates uniformly without replacement.
There are $\binom{10}{2}=45$ pairs, of which $\binom72=21$ contain no passing
candidate. Thus
\begin{equation}
 P(\hbox{at least one passes})
 =1-\frac{\binom72}{\binom{10}2}=\frac{24}{45}=\frac8{15}.
 \label{exec19:evalcode-pool}
\end{equation}
The successful count in that formula is the number among all ten candidates.
Substituting the random number of successes in the selected pair instead gives
a different quantity; exact enumeration gives average $16/135$, not $8/15$.
The source inconsistency must therefore be resolved before it becomes a
benchmark implementation or an accepted numerical comparison.
\citep{exec19_surveys_reference}

An at-least-one success event does not specify how the deployed system finds
that success. If it chooses one of the two sampled candidates uniformly,
its chance of selecting a passing candidate remains $3/10$. Reaching $8/15$
as an output-success rate requires an oracle or a suitably qualified selector
that can recognize the passing candidate whenever one is present. A selector
using public tests, a model judge or analyst review has its own errors and
costs. For causal code, even a perfect test-suite selector establishes only
success under those tests; it does not establish identification assumptions or
the relevance of the estimand.~\citep{exec19_surveys_reference}

\subsection{Retain the attempted-request denominator}
\label{exec19:evalcode-pipeline}

Consider a second authored example with twenty attempted causal coding tasks.
Sixteen produce runnable outputs; twelve pass their unit tests; seven preserve
the declared causal target and information clock; five satisfy the full output
contract, including evidence and uncertainty. These nested counts are stipulated
to illustrate measurement design, not observations from a model. The two
headline rates are different:
\begin{equation}
 r_{\rm tests\mid runnable}=\frac{12}{16}=\frac34,
 \qquad
 r_{\rm accepted\mid attempted}=\frac5{20}=\frac14.
 \label{exec19:evalcode-denominators}
\end{equation}
A report containing only the first rate would omit failed execution and
scientific-contract failures. The pipeline should retain the request identifier,
eligibility decision, every candidate and retry, selected result, failure or
abstention reason, and final adjudication. The rates need not be multiplied
under an independence assumption: the example directly counts nested outcomes.
Correlated prompts and shared graph families also mean these twenty rows do
not automatically justify a binomial interval or twenty independent units of
evidence.~\citep{exec19_surveys_reference}

This requirement follows the surveys' limitations rather than treating a
benchmark as a universal standard. The code survey stresses that execution
metrics depend on unit-test quality and omit attributes such as maintainability,
efficiency, stability and repository dependencies. Human review adds expertise
but also cost, variance and reproducibility issues. Model judges can inherit
position, verbosity, self-enhancement and reasoning biases. The evaluation
survey similarly treats evaluator variability and the validity of the testing
system as questions requiring scrutiny. A causal workflow adds its own
target, graph, clock, support and refusal semantics to this broader software
assessment.~\citep{exec19_surveys_code,exec19_surveys_evaluation}

\subsection{Count costs on all attempts and benefits on accepted outputs}
\label{exec19:evalcode-cost}

Suppose generation costs USD 2 and the test harness costs USD 1 for every
attempt, while specialist review costs USD 5 for each test-passing output.
Suppose an accepted output is worth USD 30 under an explicitly hypothetical
valuation. For the twenty-task construction,
\begin{align}
 C&=20(2+1)+12(5)=120\quad\text{USD},
 \label{exec19:evalcode-costs}\\
 C/N_{\rm accepted}&=120/5=24\quad\text{USD per accepted output},
 \label{exec19:evalcode-unitcost}\\
 B-C&=5(30)-120=30\quad\text{USD in total}.
 \label{exec19:evalcode-net}
\end{align}
Net value is USD 1.50 per attempted task. Equating all twelve test passes with
accepted outputs would instead produce USD 240 net and overstate value by
USD 210. These prices and benefits are assumptions; none is a model-provider
quote, observed analyst saving or bank ROI. The example demonstrates why the
all-attempt cost ledger and acceptance rubric must be frozen together before
a prototype trial. A real experiment should additionally include setup, retries,
timeouts, correction, dependency maintenance and the cost of rejected work.
\citep{exec19_surveys_reference}

\subsection{A public evaluation design for the proposed causal assistant}

The proposal can test these ideas without authentic bank records. Generate
declared structural and tabular examples with known interventions, ambiguous
graphs, missing outcomes, invalid clocks and unsupported requests. Separate
development graph families from evaluation families and retain original
generation seeds and source versions. Score formulation, diagram semantics,
identification, code execution, numerical output, uncertainty, reporting and
appropriate refusal separately. The classifier or judge for each stage must
itself be checked on counterexamples; agreement with a fluent explanation is
not an independent causal oracle.

Use unit tests for declared numerical and interface properties, source-grounded
checks for algorithm assumptions, and adjudicated scientific review where the
target is ambiguous. Report automatic, human and model-judge decisions with
their disagreement rather than substituting one for another invisibly.
Randomize order when testing a model judge, retain its version and prompt, and
evaluate selection separately from candidate generation. The survey discussions
motivate these dimensions but do not prove that this proposed protocol is
calibrated. Historical benchmarks can supply leads and adverse examples;
current package adoption requires its own pinned originals and conformance.
\citep{exec19_surveys_evaluation,exec19_surveys_code}

The next experiment is consequently a bounded test of whether the assistant
preserves a declared causal contract at a useful all-request cost. Public and
synthetic tasks can expose wrong denominators, selection optimism, unsafe
conversions and semantic errors before costly bank experiments. They cannot
verify an institution's actual assignment mechanism, audit truth or realized
economic return. Those remain later questions after the prototype demonstrates
why further evidence would be worth acquiring.
